\section*{Ethical Considerations}
\addcontentsline{toc}{section}{Ethical Considerations}

This work analyzes a bribery pattern that could be harmful if applied to a production system. We
therefore keep the \emph{attack} machinery local-only: the proof-of-concept runs against synthetic
committees, the \textsc{PayToEquivocate} escrow is a Python model that omits the engineering required
to weaponize a live contract, and the equivocation leg is never executed on any shared network---a
runtime gate refuses any non-local chain-id. The live measurement of \S\ref{sec:measured} is strictly
\emph{read-only}: it reads already-public ledger data (historical double-sign slashes, ordinary bridge
fills) over public endpoints, handles no funds or keys, mounts no attack, and surfaces only on-chain
events the ledger already exposes---no off-chain de-anonymization.

The benefit of publication is defensive: the model identifies exactly which design choices close the
window, and the case study and measurement show why fast evidence and finality-gated settlement
protect deployed systems. If we uncover an exploitable configuration on a named deployed system, we
will notify the affected project maintainers before public disclosure, share only the minimum details
needed for remediation during the embargo, and gate any exploit-enabling artifact components until
the relevant defenses are deployed or the maintainers approve release.
